\chapter{大規模脳ネットワークの発見と3大ネットワーク}
\epigraph{``Dūraṅgamaṁ ekacaraṁ, asarīraṁ guhāsayaṁ;\\
Ye cittaṁ saññamessanti, mokkhanti mārabandhanā.''\\
\small（遠くへ彷徨い、単独で行き、形体なく、胸の洞窟に潜むこの心。これを制御する者のみが、魔境の束縛から脱出する。）}{\textit{Dhammapada}, 37}

脳機能局在論からネットワーク論へのパラダイムシフトは、20世紀末から21世紀初頭にかけての安静時機能的磁気共鳴画像法（rs-fMRI）の発展によってもたらされた。
脳は単一領域が独立して機能を担うのではなく、遠隔に散らばる領域群が同期振動する「機能的結合ネットワーク（Large-Scale Brain Networks）」として駆動していることが明らかとなったのである。
その中核を成すのが、以下の3大ネットワークである。

\begin{itemize}
  \item \textbf{DMN（デフォルト・モード・ネットワーク / Default Mode Network）}：\\
  2001年、マーカス・レイクル（Marcus Raichle）らによって同定された。被験者が特定の認知課題を行わず「ぼんやりと安静にしている状態」で活動が高まり、逆に外部のタスクに集中すると活動が低下（タスク負相関）する領域群（内側前頭前皮質、後帯状皮質、下頭頂小葉など）。自己参照的思考、過去のエピソード記憶の想起、未来のシミュレーション、自己の物語（ナラティブ）の維持を司る。
  \item \textbf{CEN（中央実行系ネットワーク / Central Executive Network）}：\\
  背外側前頭前皮質（dlPFC）や後頭頂皮質（PPC）を中心とし、外部環境のタスク遂行、ワーキングメモリ、論理的推論、問題解決など、意識的・能動的な計算処理を実行するネットワーク。DMNとは基本的に排他・拮抗する関係にある。
  \item \textbf{SN（顕著性ネットワーク / Salience Network）}：\\
  2007年、ウィリアム・シーリー（William Seeley）やヴィノッド・メノン（Vinod Menon）らによって提唱された。前帯状皮質前部（dACC）と島皮質前部（anterior insula）を中心とする。体内外の刺激から「生存や目標にとって何が最も顕著（salient）か」を検知し、DMNとCENの活動バランスを動的に切り替えるスイッチングの役割を果たす。
\end{itemize}

\section{DMNの働き}
\label{section:DMNの働き}
DMNの本質とは、生体内部の自己防衛モジュール、すなわち「自我（Ego）」の利益（profit）と脅威（threat）のパラメータを評価し続ける自己参照型シミュレータである。その評価基準は「安全・快（profit）」と「その喪失・痛み（threat）」にあり、常にモニタリングされているのは、自身が所属しているという社会的な序列（ヒエラルキー）における、Egoと他者との相対的ポジションである。

また、DMNはEgoの経験を説明可能な「理由」を絶え間なく求め続ける。身体や情動の先行出力に対して納得できる因果関係を見出そうと、暫定的なナラティブ仮説が神経系疑似ベイズ推論機によって確定し、強固な確信へと結像するまで、環境および記憶から情報を貪欲に収集・サンプリングし続けるのである。このとき、脳の計算資源を節約し不確実性を手早く解消するため、採用される仮説は厳密な真実ではなく、直感的に「納得できる短絡的なもの」が選ばれることが多い。

脅威として確定した確信は、扁桃体を中心とする大脳辺縁系を急激に励起させ、上位の認知プロセスに情動的割り込み（インターラプション）をかけて「怒り」や「恐怖」といった生理反応を強制惹起する。このとき確定するナラティブの確信は、アトミック（素因的）な水準まで分解すると、自己と他者のヒエラルキー配置に関する極めて少数の基本パターンへと還元・整理できる。

\begin{enumerate}[label=(\arabic*)]
\item \textbf{上位確信（profit）}：「自分はヒエラルキーの上位にある」という結論（自己保存の安全・快）
\item \textbf{下位確信（threat）}：「自分はヒエラルキーの下位にある」という結論（屈服・劣等・剥奪の脅威）
\item \textbf{排除確信（threat）}：「自分はヒエラルキーの上位であり、下位の相手は不浄・汚物である」という結論（接触回避・排除）
\end{enumerate}

人間に強烈な「嫌悪」の情動が引き起こされるのは、まさにこの第3のパターンが確定した瞬間である。

人間が織りなす苦悩や葛藤、愛憎劇のドラマは、一見すると無限の深淵と複雑さを湛えているように思える。しかし、物理学における「カオス」や「フラクタル幾何学」が明かしたように、自然界における最も複雑怪奇な挙動――三体問題の予測不能な軌道や、乱流の渦、砕波のフラクタル構造――を生み出している根源は、驚くほど単純な非線形方程式の反復にすぎない。

DMNという自己参照ループもまた、これと完全に同型である。私たちが「複雑な人生の物語」と呼んでいるものは、自己防衛とヒエラルキー配置を巡るごく少数の基本アルゴリズムが、納得できる理由を求める短絡的な作話と環境ノイズを巻き込み、幾重にもフィードバックを繰り返した結果として立ち現れる、見かけ上のカオス的フラクタルにすぎないのだ。

だからこそ、私たちは人間の心という現象を必要以上に神秘化し、その表層の複雑さに眩惑されてはならない。解くべき方程式は、極めて単純に捉えるべきである。どれほど錯綜したナラティブの迷宮であっても、その深層で駆動しているのは、常にアトミックな利害判定関数と、わずかな確信パターンの収束力学にほかならないのだから。

\section{SNのスイッチング機能とDMN・CENのシーソー効果}

DMN（自己参照・過去未来の反芻・雑念）とCEN（外部タスク遂行・論理演算）は、どちらか一方が活動しているときにもう一方が抑制される「反相関（anti-correlation）」、すなわちシーソーのような拮抗関係にある。
このシーソーの傾きを能動的に制御し、DMNの暴走を沈静化させる物理的レバーとして機能するのが\textbf{顕著性ネットワーク（SN）}である。

\subsection{メノンらのトリプルネットワーク・モデル（2011年）}
スタンフォード大学のヴィノッド・メノン（Vinod Menon）とルシン・ウディン（Lucina Uddin）らは、fMRIデータと因果結合性解析（Granger因果性や動的因果モデリング）を用いた一連の実験を通じて「トリプルネットワーク・モデル」を確立した。
この実験研究が証明したのは、SN（特に右島皮質前部：rAI）が\textbf{「因果的司令塔（Causal Outflow Hub）」}として振る舞い、感覚刺激や内受容感覚の突発的な顕著性を検知した瞬間、DMNを物理的に下方抑制（ダウンレギュレーション）し、同時にCENを立ち上げるという動的スイッチングを実行している事実である。

被験者が「自己の将来の不安や後悔」といった自己言及的ナラティブ（DMN駆動状態）に陥っているとき、後帯状皮質（PCC）や内側前頭前皮質（mPFC）の血流シグナルは急増する。
しかし、被験者に「今この瞬間の呼吸の皮膚感覚」や「足の裏にかかる重力圧」といった、\textbf{直接的な身体感覚・内受容感覚（Interoception）へ注意を急峻にロックオン}させると、次の神経動態がミリ秒〜秒単位で観測される。
\begin{enumerate}[label=(\arabic*)]
  \item 島皮質前部（AI）および背側前帯状皮質（dACC）からなる\textbf{SNが急激に同期・賦活化}する。
  \item 活性化したSNが抑制性シグナルを送り出し、暴走していた\textbf{DMN（PCC/mPFC）の活動を一気に底へと押し下げる}。
  \item 自己参照的なストーリー生成が物理的に遮断され、外部受容・作業記憶を司るCEN、あるいは感覚野へのリソース再配分が行われる。
\end{enumerate}



%\chapter{大規模脳ネットワークの発見と3大ネットワーク}
%\epigraph{``Dūraṅgamaṁ ekacaraṁ, asarīraṁ guhāsayaṁ;\\
%Ye cittaṁ saññamessanti, mokkhanti mārabandhanā.''\\
%\small（遠くへ彷徨い、単独で行き、形体なく、胸の洞窟に潜むこの心。これを制御する者のみが、魔境の束縛から脱出する。）}{\textit{Dhammapada}, 37}
%脳機能局在論からネットワーク論へのパラダイムシフトは、20世紀末から21世紀初頭にかけての安静時機能的磁気共鳴画像法（rs-fMRI）の発展によってもたらされた。
%脳は単一領域が独立して機能を担うのではなく、遠隔に散らばる領域群が同期して同期振動する「機能的結合ネットワーク（Large-Scale Brain Networks）」として駆動していることが明らかとなったのである。
%その中核を成すのが、以下の3大ネットワークである。
%
%\begin{itemize}
%  \item \textbf{DMN（デフォルト・モード・ネットワーク / Default Mode Network）}：\\
%  2001年、マーカス・レイクル（Marcus Raichle）らによって同定された。被験者が特定の認知課題を行わず「ぼんやりと安静にしている状態」で活動が高まり、逆に外部のタスクに集中すると活動が低下（タスク負相関）する領域群（内側前頭前皮質、後帯状皮質、下頭頂小葉など）。自己参照的思考、過去のエピソード記憶の想起、未来のシミュレーション、自己の物語（ナラティブ）の維持を司る。
%  \item \textbf{CEN（中央実行系ネットワーク / Central Executive Network）}：\\
%  背外側前頭前皮質（dlPFC）や後頭頂皮質（PPC）を中心とし、外部環境のタスク遂行、ワーキングメモリ、論理的推論、問題解決など、意識的・能動的な計算処理を実行するネットワーク。DMNとは基本的に排他・拮抗する関係にある。
%  \item \textbf{SN（顕著性ネットワーク / Salience Network）}：\\
%  2007年、ウィリアム・シーリー（William Seeley）やヴィノッド・メノン（Vinod Menon）らによって提唱された。前帯状皮質前部（dACC）と島皮質前部（anterior insula）を中心とする。体内外の刺激から「生存や目標にとって何が最も顕著（salient）か」を検知し、DMNとCENの活動バランスを動的に切り替えるスイッチングの役割を果たす。
%\end{itemize}
%\subsection{DMNの働き}
%DMNの本質とは、生体内部の自己防衛モジュール、すなわち「自我（Ego）」の利益（profit）と脅威（threat）のパラメータを評価し続ける自己参照型シミュレータである。その評価基準は「安全・快（profit）」と「その喪失・痛み（threat）」にあり、Egoの経験を説明可能な「理由」を絶え間なく求め続ける。常にモニタリングされているのは、自身が所属しているという確信$\Omega$ のヒエラルキーにおける、Egoと他者との相対的ポジションである。また、DMNはEgoの経験を説明可能な「理由」を絶え間なく求め続け、その暫定的なナラティブ仮説が神経系疑似ベイズ推論機によって確定し、確信へと結像するまで、環境および記憶から情報を収集・サンプリングし続ける。このとき、確信は納得できる短絡的なものが選ばれることが多い。
%
%脅威として確定した確信は、扁桃体を中心とする大脳辺縁系を急激に励起させ、上位の認知プロセスに情動的割り込み（インターラプション）をかけて「怒り」や「恐怖」といった生理反応を強制惹起する。このとき確定するナラティブの確信は、アトミック（素因的）な水準まで分解すると、自己と他者のヒエラルキー配置に関する極めて少数の基本パターンへと還元・整理できる。
%
%\begin{enumerate}[label=(\arabic*)]
%\item \textbf{上位確信(profit)}：「自分はヒエラルキーの上位にある」という結論（自己保存の安全・快）
%\item \textbf{下位確信(threat)}：「自分はヒエラルキーの下位にある」という結論（屈服・劣等・剥奪の脅威）
%\item \textbf{排除確信(threat)}：「自分はヒエラルキーの上位であり、下位の相手は不浄・汚物である」という結論（接触回避・排除）
%\end{enumerate}
%
%
%人間に強烈な「嫌悪」の情動が引き起こされるのは、まさにこの第3のパターンが確定した瞬間である。
%
%人間が織りなす苦悩や葛藤、愛憎劇のドラマは、一見すると無限の深淵と複雑さを湛えているように思える。しかし、物理学における「カオス」や「フラクタル幾何学」が明かしたように、自然界における最も複雑怪奇な挙動—三体問題の予測不能な軌道や、乱流の渦、砕波のフラクタル構造—を生み出している根源は、驚くほど単純な非線形方程式の反復にすぎない。
%
%DMNという自己参照ループもまた、これと完全に同型である。私たちが「複雑な人生の物語」と呼んでいるものは、自己防衛とヒエラルキー配置を巡るごく少数の基本アルゴリズムが、環境ノイズを巻き込んで幾重にもフィードバックを繰り返した結果として立ち現れる、見かけ上のカオス的フラクタルにすぎないのだ。
%
%だからこそ、私たちは人間の心という現象を必要以上に神秘化し、その表層の複雑さに眩惑されてはならない。解くべき方程式は、極めて単純に捉えるべきである。どれほど錯綜したナラティブの迷宮であっても、その深層で駆動しているのは、常にアトミックな利害判定関数と、わずかな確信パターンの収束力学にほかならないのだから。
%

%\subsection{DMNの働き}
%DMNの本質とは、生体内部の自己防衛モジュール、すなわち「自我（Ego）」の利益（profit）と脅威（threat）のパラメータを評価し続ける自己参照型シミュレータである。その評価基準は「安全・快（profit）」と「その喪失・痛み（threat）」にあり、常にモニタリングされているのは、自身が所属しているという確信$\Omega$ のヒエラルキーにおける、Egoと他者との相対的ポジションである。また、DMNはEgoの経験を説明可能な「理由」を絶え間なく求め続け、その暫定的なナラティブ仮説 $\hat{\Omega}$ が神経系疑似ベイズ推論機によって確定し、確信 $\Omega$ へと結像するまで、環境および記憶から情報を収集・サンプリングし続ける。
%
%脅威として確定した確信 $\Omega_{\mathrm{threat}}$ は、扁桃体を中心とする大脳辺縁系を急激に励起させ、上位の認知プロセスに情動的割り込み（インターラプション）をかけて「怒り」や「恐怖」といった生理反応（$\Lambda_{\mathrm{anger}}, \Lambda_{\mathrm{fear}}$）を強制惹起する。このとき確定するナラティブの確信 $\Omega$ は、アトミック（素因的）な水準まで分解すると、自己と他者のヒエラルキー配置に関する極めて少数の基本パターンへと還元・整理できる。\begin{enumerate}[label=(\arabic*)]\item \textbf{上位確信}：「自分はヒエラルキーの上位にある」という結論（自己保存の安全・快）\item \textbf{下位確信（$\Omega_{\mathrm{threat}}$）}：「自分はヒエラルキーの下位にある」という結論（屈服・劣等・剥奪の脅威）\item \textbf{排除確信}：「自分はヒエラルキーの上位であり、下位の相手は不浄・汚物である」という結論（接触回避・排除）\end{enumerate}人間に強烈な「嫌悪（$\Lambda_{\mathrm{disgust}}$）」の情動が引き起こされるのは、まさにこの第3のパターンが確定した瞬間である。
%

